\chapter{Evaluation Methodology and Experimental Design}
\label{sec:evaluation}

This chapter should explain how the thesis evaluates the system without overstating what the current data can support. The repository already contains a concrete binary-analysis harness, and the final chapter should align tightly with that implementation.

\section{Evaluation Goals}

\begin{itemize}
\item Measure whether tool-orchestrated workflows outperform weaker baselines on binary-analysis tasks.
\item Separate the effects of tool access, workflow topology, prompt shaping, and review strictness.
\item Capture both successful analyses and structured failure modes such as analysis errors, validator blocks, and judge failures.
\end{itemize}

\begin{itemize}
\item \TODO{State the primary evaluation objective in thesis-ready language.}
\item \TODO{List the secondary ablation questions and map each one to a configuration family.}
\end{itemize}

\section{Benchmark Corpus and Task Model}

The maintained binary benchmark uses an experimental corpus of authored executables and executable variants. The evaluation unit is a sample-task pair rather than only a sample. This detail is important because it lets the thesis distinguish broad understanding tasks from focused recovery tasks.

\begin{itemize}
\item Sample variants cover conditions such as stripping, packing, deceptive strings, anti-analysis behavior, and configuration decoding.
\item Tasks include default analysis as well as focused recovery questions such as config extraction, decode-flow recovery, packing triage, and behavioral reconstruction.
\item The corpus supports both breadth-oriented sweeps and depth-oriented studies on selected executables such as \texttt{config\_decoder\_test}.
\end{itemize}

\begin{itemize}
\item \TODO{Describe the experimental corpus, including how many source programs, executable variants, and sample-task pairs were used for the final thesis experiments.}
\item \TODO{Add a short table summarizing representative samples and the reasoning skills they are intended to stress.}
\item \TODO{Explain why \texttt{config\_decoder\_test} became a centerpiece for deeper analysis.}
\end{itemize}

\section{Single Runs, Sweeps, and Repetitions}

The harness supports both single-run evaluation and experiment sweeps. Maintained sweeps are baseline-first and primarily one-variable-at-a-time. The thesis should explain this choice clearly, because it affects how causal claims can be interpreted.

\begin{itemize}
\item A global baseline defines the default workflow configuration.
\item Variant families change one configuration surface at a time.
\item Some families insert local family baselines where a fair comparison requires it.
\item Repetitions support stronger claims about stability and significance.
\end{itemize}

\begin{itemize}
\item \TODO{Explain why the broad sweep uses one-variable-at-a-time comparisons instead of a full factorial design.}
\item \TODO{Document the final repetition counts used for the headline experiments and for the decoder depth study.}
\item \TODO{Clarify how repaired experiments are handled when failed tasks are rerun and merged.}
\end{itemize}

\section{Baselines and Experimental Variables}

The current harness exposes several variable families that can be studied empirically:

\begin{itemize}
\item query verbosity
\item worker subagent profile
\item worker prompt shape and worker role-prompt mode
\item tool availability
\item architecture preset
\item pipeline preset
\item validator topology
\item validator review level
\end{itemize}

\begin{itemize}
\item \TODO{Decide which of these families will be discussed in the thesis body versus moved to appendix material.}
\item \TODO{Explain the rationale for the chosen baseline in the final experiments.}
\item \TODO{Document the deeper \texttt{config\_decoder\_test} study separately from the broad sweep.}
\end{itemize}

\section{Judging and Scoring}

Run outputs are judged with a structured binary-analysis rubric and converted into continuous scores. This thesis should make the judging methodology explicit, including both its usefulness and its limitations.

\begin{itemize}
\item Per-task judge outputs preserve structured feedback and an overall score.
\item Legacy pass or fail thresholds still exist, but sweep analysis should focus on score deltas, success rates, and error rates.
\item The harness preserves analysis outcomes such as completed, analysis error, no result, judge error, and validator blocked.
\end{itemize}

\begin{itemize}
\item \TODO{Describe the rubric dimensions used by the judge and justify why they align with reverse-engineering quality.}
\item \TODO{Explain whether human spot-checking was used to validate judge outputs.}
\item \TODO{Discuss the difference between completed-but-weak analyses and infrastructure failures.}
\end{itemize}

\section{Statistical Treatment and Comparison Strategy}

The final thesis should not report only descriptive movement. The harness supports paired comparisons and limited significance analysis when repetitions or enough matched task pairs exist.

\begin{itemize}
\item Replicate-level comparison when repeated runs exist.
\item Paired-task comparison across matched sample-task pairs.
\item Reporting of score deltas, task success deltas, and significance outputs where supported.
\end{itemize}

\begin{itemize}
\item \TODO{State the statistical tests actually used in the final thesis and justify their assumptions.}
\item \TODO{Document how missing or recovered tasks are treated in the analysis set.}
\item \TODO{Explain what level of evidence is sufficient for claims in the results chapter.}
\end{itemize}
